\documentclass[12pt,anyside]{book}

% ============================================================
% PREAMBLE
% ============================================================
\usepackage[utf8]{inputenc}
\usepackage[T1]{fontenc}
\usepackage{lmodern}
\usepackage{microtype}
\usepackage[a4paper,margin=2.5cm]{geometry}
\usepackage{setspace}
\onehalfspacing
\usepackage{amsmath,amssymb,amsthm}
\usepackage{mathtools}
\usepackage{booktabs}
\usepackage{tabularx}
\usepackage{comment}
\usepackage{graphicx}
\usepackage{float}
\usepackage{caption}
\usepackage{subcaption}
\usepackage{tikz}
\usepackage{pgfplots}
\pgfplotsset{compat=1.18}
\usepackage{hyperref}
\hypersetup{
    colorlinks=true,
    linkcolor=blue,
    citecolor=blue,
    urlcolor=blue
}
\usepackage{cleveref}
\usepackage{enumitem}
\usepackage{fancyhdr}
\usepackage{titlesec}
\usepackage{tocloft}
\usepackage{appendix}

% Theorem environments
\theoremstyle{definition}
\newtheorem{definition}{Definition}[chapter]
\newtheorem{assumption}{Assumption}[chapter]
\theoremstyle{plain}
\newtheorem{proposition}{Proposition}[chapter]
\newtheorem{theorem}{Theorem}[chapter]
\newtheorem{corollary}{Corollary}[chapter]
\theoremstyle{remark}
\newtheorem{remark}{Remark}[chapter]

% Header / Footer
\pagestyle{fancy}
\fancyhf{}
\fancyhead[LE]{\leftmark}
\fancyhead[RO]{\rightmark}
\fancyfoot[C]{\thepage}
\renewcommand{\headrulewidth}{0.4pt}

% Title formatting
\titleformat{\chapter}[display]
{\normalfont\huge\bfseries}{\chaptertitlename\ \thechapter}{20pt}{\Huge}
\titleformat{\section}{\normalfont\Large\bfseries}{\thesection}{1em}{}
\titleformat{\subsection}{\normalfont\large\bfseries}{\thesubsection}{1em}{}

% ============================================================
% DOCUMENT
% ============================================================
\begin{document}

% ----------------------------------------------------------
% FRONT MATTER
% ----------------------------------------------------------
\frontmatter

\begin{titlepage}
\centering
\vspace*{2cm}
{\Huge\bfseries Equity-Share Corporatism\\[0.4cm] with Frozen Prices\par}
\vspace{1.5cm}
{\LARGE A New Institutional Architecture for\\[0.3cm] National Greatness\par}
\vspace{2cm}
{\Large Samir Amier Saliem Boulos\par}
\vspace{1cm}
{\large A Vision for Australia\par}
\vfill
{\large 2026\par}
\end{titlepage}

\chapter*{Dedication}
\thispagestyle{empty}
To the people of Australia, and to every citizen who believes that a nation can be both prosperous and just, both free and purposeful.
\vspace{2cm}
\begin{center}
\textit{``The purpose of an economy is not merely to allocate scarce resources,\\ but to enlarge the space in which human flourishing becomes possible.''}
\end{center}

\chapter*{Preface}
\addcontentsline{toc}{chapter}{Preface}

This book is written for economists, political scientists, public servants, parliamentarians, and every Australian who senses that the present institutional settlement, while historically successful, is no longer sufficient to secure the highest possible trajectory of national flourishing.

It presents a complete, internally consistent economic and political architecture. The architecture is deliberately constrained: a flat personal income tax of ten per cent; universal freezing of official commodity prices; a permanent thirty-per-cent government equity claim in every corporation (partnerships and sole traders excepted); unrestricted sales-contingent wage adjustment; full transparency of bank interest rates; and the deliberate insulation of official prices from the forces of supply and demand. These constraints are not rhetorical flourishes. They are binding design parameters. Every subsequent mechanism, every second-best corrective, and every comparative claim is developed strictly inside that feasible set.

The analysis is rigorous where rigour is required. Mathematical statements of equilibrium, incentive compatibility, misallocation bounds, and growth accounting appear where they authenticate the argument. Yet the exposition is deliberately accessible. Technical sections are preceded by plain-language summaries so that a senior official, a parliamentary adviser, or an informed citizen can grasp the logic without mastering every derivation. The mathematics is present to discipline the claims, not to exclude the reader.

Australia stands at a hinge moment. Productivity growth has slowed. Housing and essential goods remain under persistent pressure. The tax-transfer system has become complex. Public ownership is limited and ad hoc. This book offers a coherent alternative: an equity-share corporatist quantity-constrained mixed economy capable, if implemented with discipline, of elevating Australia to a position of unambiguous leadership among nations.

The chapters that follow consolidate the entire body of analysis developed to date. They are written as a single, continuous argument.

\begin{comment}
NOTES FOR PREFACE:
- Emphasise accessibility + rigour balance.
- State clearly that all seven original constraints remain inviolable.
- Position the book as both scholarly contribution and practical blueprint for government.
- Tone: confident, non-partisan, national-purpose oriented.
\end{comment}

\newpage

\tableofcontents
\listoffigures
\listoftables

% ----------------------------------------------------------
% MAIN MATTER
% ----------------------------------------------------------
\mainmatter

\chapter{The Problem of Institutional Design}
\label{ch:problem}

Institutions are not ornaments of economic life. They are the rules that determine which actions are rewarded, which risks are borne, and which futures become feasible. When the rules align private incentives with social value, societies accumulate capital, knowledge and trust. When the rules systematically misalign those incentives, even abundant resources and talented people produce mediocrity. This chapter establishes why institutional architecture is the decisive variable in long-run national performance, diagnoses the principal limitations of the contemporary Australian settlement, states with precision the seven binding constraints that define the alternative architecture developed in this book, and maps the logical path the subsequent chapters will follow.

\section{Why Institutions Matter}

At the most fundamental level, an economy is a system for coordinating the specialised activities of millions of people who cannot possibly know one another’s circumstances in detail. Coordination can occur through three pure mechanisms: voluntary exchange guided by relative prices, hierarchical command, or repeated bargaining inside stable organisations. Every real economy is a mixture of the three. The particular mixture chosen—by design or by historical accident—determines the information that is generated, the incentives that agents face, and therefore the rate at which living standards can rise.

Relative prices perform a unique cognitive function. When a price is free to move, it aggregates dispersed knowledge about scarcity and preference into a single, publicly observable statistic. Agents need not know why copper has become scarcer; they need only observe that its price has risen and adjust their plans accordingly. Residual claimancy—the right to the net surplus after all contractual obligations are met—performs a complementary incentive function. The residual claimant is the residual risk-bearer and therefore the party with the strongest interest in maximising the joint surplus. When residual claims are secure and concentrated, effort, innovation and careful stewardship of assets are rewarded. When residual claims are attenuated or politicised, those same behaviours become less attractive.

These two institutions—flexible relative prices and well-defined residual claims—do not exhaust the list of determinants of long-run performance, but they are foundational. Douglass North’s work on the institutional foundations of economic growth, Daron Acemoglu and James Robinson’s analysis of inclusive versus extractive institutions, and Oliver Williamson’s transaction-cost theory of governance all converge on the same conclusion: societies that protect residual claims and allow prices to convey information tend to generate sustained increases in productivity; societies that systematically weaken residual claims or suppress price signals tend to stagnate or decline relative to their potential.

The implication for national strategy is direct. Australia’s future prosperity will be shaped less by the discovery of additional natural resources or by the temporary stimulus of public expenditure than by the quality of the rules that govern ownership, pricing, taxation and coordination. A coherent redesign of those rules can shift the entire growth–equity frontier outward. An incoherent or purely incremental approach cannot.

\section{The Limits of the Contemporary Australian Settlement}

The Australian institutional settlement that emerged from the reforms of the 1980s and 1990s remains, by global standards, a high-performing mixed economy. Markets allocate the great majority of goods and services. Private residual claims predominate outside a limited set of government business enterprises. An independent central bank targets low and stable inflation. A progressive personal income tax and a comprehensive transfer system provide a substantial degree of redistribution. Compulsory superannuation has created a large pool of patient capital. These arrangements have delivered decades of uninterrupted growth and high average living standards.

Yet the same settlement now exhibits clear and measurable limitations.

First, the personal income-tax schedule remains steeply progressive. For the 2026–27 income year the rates for resident individuals are zero up to \$18\,200, 15 per cent between \$18\,201 and \$45\,000, 30 per cent between \$45\,001 and \$135\,000, 37 per cent between \$135\,001 and \$190\,000, and 45 per cent above \$190\,000, together with a 2 per cent Medicare levy. Company tax is dual-rate: 25 per cent for base-rate entities and 30 per cent otherwise. The Goods and Services Tax stands at 10 per cent. The resulting system is complex, generates significant compliance costs, and creates high effective marginal rates for many households once interactions with the transfer system are taken into account.

Second, public ownership is narrow and ad hoc. The Commonwealth maintains a small number of government business enterprises—NBN Co, Australia Post, Snowy Hydro, Defence Housing Australia and a handful of others—while the states retain ownership of certain energy, water and transport assets. There is no systematic residual claim by the public on the profits of the corporate sector as a whole. Government spending has risen to approximately 28–29 per cent of GDP, yet this expansion has occurred primarily through transfers and consumption rather than through a coherent equity stake in productive enterprise.

Third, relative prices remain the dominant allocative mechanism, subject only to limited sectoral overlays. The recent introduction of an “excessive-pricing” prohibition applicable to supermarket chains with annual revenue exceeding \$30 billion represents a modest departure from pure market determination, but the great bulk of prices continue to clear markets. The Reserve Bank of Australia retains an independent cash-rate instrument whose primary objective is an inflation target of 2–3 per cent on average over the cycle. Banks set contractual interest rates freely, subject to prudential regulation and increasing disclosure requirements.

Fourth, multi-factor and labour productivity growth have been weak for an extended period. Official data show labour productivity growth near zero or only marginally positive in recent years, while aggregate GDP growth has been supported to a significant degree by population increase. Real wages have accordingly stagnated or declined in certain episodes. These outcomes are not the result of a single policy failure; they reflect the cumulative effect of regulatory complexity, infrastructure bottlenecks, and an incentive structure that does not sufficiently reward measured productivity gains at the firm level.

Taken together, these features describe a high-income mixed economy that has lost institutional momentum. Incremental reform inside the existing paradigm can produce marginal improvements. It cannot, however, generate the discontinuous shift in national performance that a deliberately redesigned architecture can achieve.

\section{The Seven Binding Constraints}

The alternative architecture developed in this book is defined by seven constraints that are treated as inviolable. They are not bargaining positions or aspirational goals; they are the fixed parameters inside which every subsequent mechanism, corrective and comparative claim must be constructed.

\begin{enumerate}[label=\arabic*)]
\item A flat personal income tax of 10 per cent applied to all taxable income of natural persons.
\item The official price of every commodity and product is frozen at a specific, government-determined level and remains fixed.
\item The Commonwealth holds a static 30 per cent equity claim in every corporation, whether public or private and including banks; partnerships and sole-trader structures are exempt.
\item Nominal salaries and wages may rise freely as a non-decreasing function of the firm’s sales revenue or any other performance metric chosen by the firm.
\item The preceding rules are intended to strengthen incentives for productivity and superior management.
\item Because official prices are frozen, movements in supply and demand do not affect those prices.
\item Financing institutions remain free to set contractual interest rates, provided every charge is fully transparent and contemporaneously verifiable; hidden or opaque fees are prohibited.
\end{enumerate}

These seven constraints jointly imply a distinctive institutional regime. Relative prices are removed as the primary signalling mechanism. Residual claims are shared in fixed proportion between private owners and the state. Factor prices (wages and interest rates) retain substantial flexibility. Taxation of personal income is radically simplified. The regime is therefore neither classical socialism (private residual claims remain substantial), nor classical capitalism (prices are frozen and the state is a universal co-owner), nor ordinary social democracy (progressivity is abolished and price controls are comprehensive). It is a hybrid whose precise character will be located on the spectrum of economic governance in Chapter~\ref{ch:taxonomy}.

Because the constraints are binding, the analysis that follows is an exercise in constrained optimisation. We do not ask whether a different set of constraints would be preferable. We ask how the largest possible gains in productivity, growth and social cohesion can be extracted while remaining strictly inside the feasible set defined by the seven rules.

\section{Roadmap of the Argument}

The remainder of the book proceeds in six tightly linked stages.

Chapter~\ref{ch:model} constructs a formal economic model that incorporates the seven constraints as institutional primitives. It characterises the resulting quantity-constrained equilibrium, derives the managerial incentive problem under partial public ownership, and shows how frozen prices distort capital accumulation.

Chapter~\ref{ch:solutions} develops the complete set of second-best corrective instruments that remain feasible inside the constraints. Priority rationing keyed to measured productivity, high-powered contracts financed by the private residual claim, transparent priority lending, and dividend-based fiscal stabilisation are formalised, compared and integrated into a single coherent mechanism.

Chapter~\ref{ch:comparison} places the constrained regime beside the contemporary Australian settlement using the most recent publicly available data on tax schedules, ownership structures, price-formation rules and macroeconomic outcomes. Side-by-side quantitative comparisons clarify the magnitude and direction of the institutional shift.

Chapter~\ref{ch:taxonomy} locates the regime on the multi-dimensional spectrum of economic governance. Ideal types—laissez-faire capitalism, mixed social democracy, market socialism, command socialism and fascist corporatism—are mapped, and the constrained model is shown to occupy a previously sparsely populated region best described as equity-share corporatism with frozen prices.

Chapter~\ref{ch:implementation} translates the architecture into a practical pathway for legislative and administrative action. Sequencing, agency roles, transitional safeguards and political-economy risks are examined with the needs of government officials in mind.

Chapter~\ref{ch:vision} returns to the national purpose that motivates the entire project. It argues that disciplined application of the constrained architecture can shift Australia onto a higher growth–equity frontier and thereby establish the country as an exemplar of purposeful institutional design.

The argument is cumulative. Each chapter presupposes the results of those that precede it. The mathematics authenticates the claims; the exposition ensures that the claims remain intelligible to the economists, political scientists and public officials who must ultimately judge and, if persuaded, implement them.

The choice before Australia is not between perfection and the status quo. It is between an incremental continuation of a settlement that has lost momentum and a coherent redesign that accepts clear constraints in order to unlock greater national capability. The chapters that follow are offered as a contribution to that choice.

\begin{comment}
NOTES FOR CHAPTER 1 (expanded into text above):
- First-principles discussion of relative prices, residual claims and coordination mechanisms draws on North, Acemoglu-Robinson and Williamson without requiring the reader to have prior familiarity.
- Contemporary Australian performance summarised with specific 2026-27 ATO rates, approximate government-spending share, limited SOE list, recent anti-gouging measure, RBA independence and weak productivity data.
- Seven constraints stated verbatim and declared inviolable.
- Roadmap explicitly links modelling $\to$ second-best correctives $\to$ Australian comparison $\to$ taxonomic placement $\to$ implementation $\to$ national-greatness vision.
- Language kept accessible; mathematical notation introduced only where it prepares the ground for later chapters; no redundant restatement of material that will be developed fully downstream.
\end{comment}


\chapter{A Formal Model of the Constrained Economy}
\label{ch:model}

This chapter constructs a complete formal representation of an economy that operates under the seven binding constraints stated in Chapter~\ref{ch:problem}. The purpose is not to produce an abstract theoretical curiosity, but to obtain precise, testable implications about shortages, managerial behaviour, capital accumulation and cross-market allocation. Only when those implications are derived with clarity can the second-best correctives of Chapter~\ref{ch:solutions} be designed with confidence.

The exposition proceeds in four steps. First, the economic primitives—households, firms, technology and preferences—are defined and the seven constraints are translated into mathematical statements. Second, the classical Walrasian equilibrium concept is replaced by a quantity-constrained equilibrium that is consistent with frozen official prices. Third, the managerial agency problem that arises from the mandatory thirty-per-cent public equity claim is characterised. Fourth, the implications for capital accumulation are derived from an Euler equation evaluated at distorted relative prices. The chapter closes with a concise catalogue of the inefficiencies that the subsequent chapter must address.

\section{Primitives and Institutional Constraints}

Consider a closed economy populated by a continuum of households indexed by $h\in[0,1]$ and a finite collection of firms indexed by $j\in J$. Time is discrete and infinite, $t=0,1,2,\ldots$. There are $n$ commodities.

Each household $h$ is endowed with a time endowment normalised to unity and maximises the intertemporal utility function
\begin{equation}
U^h=\mathbb{E}_0\sum_{t=0}^\infty\beta^t\,u\bigl(c_t^h,\ell_t^h\bigr),
\end{equation}
where $c_t^h\in\mathbb{R}_+^n$ is the consumption vector, $\ell_t^h\in[0,1]$ is leisure, $\beta\in(0,1)$ is the discount factor, and $u$ is strictly increasing, strictly concave and continuously differentiable. Households supply labour and own shares in firms.

Firms are partitioned into two legal categories. Let $J^C\subset J$ denote the set of corporations and $J^P=J\setminus J^C$ the set of partnerships and sole traders. Only corporations are subject to the mandatory public equity claim. The production technology of firm $j$ is given by
\begin{equation}
y_{j,t}=F_j\bigl(k_{j,t},l_{j,t},e_{j,t};A_{j,t}\bigr),
\end{equation}
where $k_{j,t}$ is the capital stock, $l_{j,t}$ is labour input, $e_{j,t}\ge0$ is managerial effort, and $A_{j,t}$ is a productivity index that may evolve exogenously or endogenously. The function $F_j$ is concave, continuous, strictly increasing in each argument, and satisfies the usual Inada conditions.

The seven binding constraints are formalised as follows.

\begin{enumerate}[label=\arabic*)]
\item Personal income tax is linear at rate $\tau=0.10$. Disposable income of household $h$ is therefore $(1-\tau)$ times the sum of labour earnings, interest receipts and distributed profits.
\item Official commodity prices are frozen: there exists a fixed vector $\bar{p}\in\mathbb{R}_{++}^n$ such that $p_{i,t}=\bar{p}_i$ for every commodity $i$ and every date $t$.
\item For every corporation $j\in J^C$ the government holds a static equity claim of $\alpha=0.30$. The private residual claim is therefore $1-\alpha=0.70$. Partnerships and sole traders ($j\in J^P$) are exempt.
\item Nominal wages $w_{j,t}$ may be adjusted freely by the firm provided the adjustment is a non-decreasing function of realised sales revenue $\bar{p}\cdot y_{j,t}$ (or of any other verifiable performance metric chosen by the firm).
\item The institutional design is intended to strengthen incentives for productivity and superior management; this intention is operationalised through the instruments derived in Chapter~\ref{ch:solutions}.
\item Because official prices are fixed, movements in supply or demand cannot alter $\bar{p}$.
\item Banks and other financing institutions set contractual interest rates freely, subject only to the requirement that every fee and charge is fully transparent and contemporaneously verifiable.
\end{enumerate}

Government revenue consists of the flat-tax receipts plus the stream of dividends accruing to the $\alpha=0.30$ equity claim. Public expenditure and any residual fiscal imbalance are left unspecified at this stage; they will be revisited when fiscal sustainability is examined.

\section{Quantity-Constrained Equilibrium}

When official prices are frozen, the classical market-clearing condition $Q_i^d(p)=Q_i^s(p)$ cannot be imposed. An alternative equilibrium concept is required.

Define the notional (unconstrained) demand and supply of each household and firm in the usual way. Because the official price vector $\bar{p}$ is fixed, these notional plans are generically inconsistent. Effective transactions are therefore determined by a rationing scheme. Let $R_t^h\in\mathbb{R}_+^n$ denote the vector of quantity constraints faced by household $h$ and let $S_t^j\in\mathbb{R}_+^n$ denote the vector of quantity constraints faced by firm $j$. Effective demand and supply are the solutions to the constrained optimisation problems
\begin{align}
\tilde{c}_t^h&=\arg\max_{c\le R_t^h}u(c,\ell)\quad\text{subject to the household budget},\\
\tilde{y}_t^j&=\arg\max_{y\le S_t^j}\pi_j(\bar{p},y,e)\quad\text{subject to the technology}.
\end{align}
A quantity-constrained equilibrium is a collection of effective plans, rationing vectors and a frozen price vector such that the aggregate feasibility condition
\begin{equation}
\int_0^1\tilde{c}_{i,t}^h\,dh\le\sum_{j\in J}\tilde{y}_{i,t}^j+\text{inventories}_{i,t}
\end{equation}
holds for every commodity $i$, and the rationing scheme is consistent with the observed transactions.

Two properties of this equilibrium concept are immediate. First, if $\bar{p}_i$ lies below the Walrasian equilibrium price $p_i^*$, then excess demand $\Delta_{i,t}=Q_{i,t}^d(\bar{p}_i)-Q_{i,t}^s(\bar{p}_i)>0$ must be eliminated by rationing. Second, once arbitrage incentives are extinguished by the price freeze, suppliers become indifferent across destinations. As shown by Albrecht, Tabarrok and Whitmeyer (2026), arbitrarily small cost differentials can then tip the entire allocation into extreme corner solutions: some markets are fully served while others are starved. This cross-market misallocation constitutes a distinct welfare loss over and above the classical Harberger triangle associated with the aggregate quantity shortfall.

The precise form of the rationing rule is left open at this stage. Chapter~\ref{ch:solutions} will select a productivity-priority rule that maximises residual social surplus inside the constraints.

\section{Managerial Incentives under Partial Public Ownership}

Consider a corporation $j\in J^C$. The manager chooses effort $e$ to maximise
\begin{equation}
V(e)=(1-\alpha)\pi_j(\bar{p},e)+b(e)-\psi(e),
\end{equation}
where $\pi_j$ is accounting profit evaluated at the frozen price vector, $b(e)$ captures private benefits of control (perk consumption, empire-building, quiet life), and $\psi(e)$ is the personal cost of effort. Both $b$ and $\psi$ are twice continuously differentiable, with $b'>0$, $b''\le0$, $\psi'>0$ and $\psi''>0$.

The first-order condition for an interior optimum is
\begin{equation}
(1-\alpha)\frac{\partial\pi_j}{\partial e}+b'(e)=\psi'(e).
\end{equation}
The corresponding social first-order condition (which internalises the government’s residual claim) is
\begin{equation}
\frac{\partial\pi_j}{\partial e}=\psi'(e).
\end{equation}
Because $\alpha=0.30>0$ and $b'(e)>0$, the manager’s privately optimal effort lies strictly below the social optimum. This under-provision of effort is the first fundamental agency cost of the mandatory public equity claim.

A second agency cost arises from the soft-budget-constraint problem. Because the government is a residual co-owner, it has a political incentive to inject additional funds when $\pi_j<0$. Rational managers anticipate this possibility and therefore internalise a lower downside risk. The result is excess investment in low-return projects and delayed exit from unproductive activities—the classic syndrome identified by Kornai.

Both agency costs will be addressed in Chapter~\ref{ch:solutions} by instruments that operate exclusively on the private residual claim and therefore leave the statutory equity share $\alpha=0.30$ untouched.

\section{Capital Accumulation and the Distorted Euler Equation}

Capital accumulation obeys the standard law of motion
\begin{equation}
k_{j,t+1}=(1-\delta)k_{j,t}+I_{j,t},
\end{equation}
where $\delta\in(0,1)$ is the depreciation rate. The representative household’s Euler equation for capital, evaluated at the frozen price vector, reads
\begin{equation}
1=\beta\mathbb{E}_t\Bigl[(1-\delta)+(1-\tau)r_{t+1}^k(\bar{p})\Bigr],
\end{equation}
where $r_{t+1}^k(\bar{p})$ is the rental rate of capital measured at official prices. In an efficient economy the same Euler equation would be evaluated at the vector of shadow prices $\lambda$ that equate marginal rates of substitution and transformation. Because $\bar{p}\neq\lambda$ in general, the capital stock converges to a steady-state level that differs from the golden-rule (or modified golden-rule) level.

Two distortions are particularly important. First, sectors whose official prices lie furthest below their shadow prices attract too little capital. Second, the overall return to capital is compressed whenever the price freeze reduces measured profits, lowering the steady-state capital intensity of the entire economy. Both effects reduce long-run consumption possibilities relative to a flexible-price benchmark.

\section{Summary of Theoretical Predictions}

The formal model yields five sharp predictions that must be confronted by any feasible corrective architecture.

\begin{enumerate}[label=\arabic*)]
\item Persistent excess demand (shortages) in every market whose frozen price lies below the would-be Walrasian price, together with the possibility of extreme cross-market misallocation once arbitrage incentives are removed.
\item Systematic under-provision of managerial effort and a tendency toward soft budget constraints arising from the mandatory public equity claim of 30 per cent.
\item Distorted sectoral allocation of capital and a lower aggregate capital intensity because the Euler equation is evaluated at official rather than shadow prices.
\item A fiscal revenue base that depends jointly on the flat 10 per cent personal tax and the dividend stream accruing to the public equity claim; compression of profits therefore threatens fiscal sustainability.
\item Incentives for agents to divert activity into black markets or informal arrangements whenever the expected return to formality is depressed by rationing and monitoring costs.
\end{enumerate}

These five inefficiencies define the design problem solved in the next chapter. Every instrument proposed there must remain strictly inside the seven binding constraints while reducing, as far as second-best considerations permit, the welfare losses identified above.

\begin{comment}
NOTES FOR CHAPTER 2 (fully expanded into the text above):
- Households, firms (corporations vs.\ partnerships/sole traders), technology $F_j(k,l,e;A)$ and preferences defined at the outset.
- Seven constraints translated into precise mathematical statements.
- Walrasian equilibrium replaced by quantity-constrained equilibrium with explicit rationing vectors; Chaos Theorem of Albrecht–Tabarrok–Whitmeyer (2026) invoked for cross-market misallocation.
- Manager’s first-order condition under $\alpha=0.30$ derived and shown to imply under-provision of effort relative to the social optimum; soft-budget-constraint problem noted.
- Euler equation written at frozen prices $\bar{p}$ and shown to generate dynamic inefficiency.
- Chapter closes with an enumerated list of the five predicted inefficiencies that Chapter~\ref{ch:solutions} must address.
- Plain-language paragraphs precede each formal block; notation is introduced only as needed; no material that belongs in later chapters is anticipated beyond the necessary forward references.
\end{comment}

\chapter{Second-Best Correctives within the Constraints}
\label{ch:solutions}

The five inefficiencies derived in Chapter~\ref{ch:model} are not optional features of the constrained regime; they are logical consequences of the seven binding rules. Because those rules cannot be relaxed, every corrective instrument must operate strictly inside the feasible set they define. This chapter constructs, compares and integrates the complete set of second-best instruments that satisfy that requirement. The objective is to minimise the residual efficiency losses while leaving the flat tax, the price freeze, the thirty-per-cent equity claim, the wage-flexibility rule, the insulation of official prices from supply and demand, and the transparency mandate for interest rates entirely untouched.

The analysis proceeds domain by domain. For each of the five predicted inefficiencies we first identify the instruments that remain feasible, then formalise their incentive and information properties, and finally select the welfare-dominant combination. The chapter concludes by assembling the selected instruments into a single, internally consistent mechanism that is incentive-compatible for private residual claimants and self-financing from the mandated public equity claim.

\section{Addressing Shortages and Misallocation}

Frozen official prices generically produce excess demand. When $\bar{p}_i<p_i^*$, the shortfall $\Delta_i=Q_i^d(\bar{p}_i)-Q_i^s(\bar{p}_i)$ must be eliminated by a non-price rule. In addition, the destruction of arbitrage incentives opens the door to the extreme cross-market misallocation characterised by the Chaos Theorem of Albrecht, Tabarrok and Whitmeyer (2026). Three instruments remain feasible inside the constraints.

\begin{description}
\item[Instrument A1 -- Productivity-priority rationing.]  
The government, exercising the voting rights that accompany its $\alpha=0.30$ equity claim, imposes a rationing protocol that ranks claimants by a verifiable productivity score $s_j$. The score may be constructed from value-added per hour, a Solow residual, or the sales-growth metric already privileged by constraint~4. The formal allocation rule is
\begin{equation}
R_i^j=\frac{s_j^\gamma}{\sum_k s_k^\gamma}\cdot Q_i^s(\bar{p}),\qquad\gamma\ge1.
\end{equation}
As $\gamma\to\infty$ the rule converges to a pure priority ordering: higher-scoring firms and households are fully served before lower-scoring claimants receive any allocation.

\item[Instrument A2 -- Non-cash tradable ration tickets.]  
The government issues a quantity of tickets equal to planned aggregate supply and distributes them initially according to the same productivity score. Tickets may be transferred among registered agents, but only at a zero official price; any secondary cash market is prohibited and monitored through the transparent banking system required by constraint~7.

\item[Instrument A3 -- Mandatory output floors.]  
Using its board representation the government votes binding production quotas $y_j\ge\underline{y}_j(\bar{p})$ calibrated to meet a target fraction of measured demand. Any resulting losses are absorbed pro-rata by the private residual claimants.
\end{description}

Instrument A2 generates strong incentives for ticket holders to reveal willingness to pay, yet the prohibition on cash side-payments is informationally demanding and creates an obvious black-market opportunity. Instrument A3 restores physical quantities but weakens cost-minimisation incentives and risks converting the soft-budget-constraint problem into a formal entitlement. Instrument A1 dominates on both incentive and information grounds: it utilises the ownership channel already mandated by constraint~3, relies on performance metrics already incentivised by constraint~4, and concentrates scarce goods on high-productivity uses. The growth-maximising parameterisation is $\gamma\to\infty$ together with periodic, transparent recalibration of the score to limit ratchet effects.

\section{Restoring Managerial Effort}

The mandatory public equity claim attenuates the manager’s residual return and thereby produces under-provision of effort. Three instruments remain feasible.

\begin{description}
\item[Instrument B1 -- High-powered contracts financed solely by the private residual claim.]  
Constraint~4 already permits wages and bonuses to rise with sales revenue. The private residual claimants (who hold 70~per cent of the equity) are therefore free to write compensation contracts of the form
\begin{equation}
w(e)=w_0+\beta\cdot(\bar{p}\cdot y(e)),\qquad\beta>0.
\end{equation}
Because the public claim is static and non-dilutable, the entire marginal return to stronger incentives accrues to the private owners who design the contract.

\item[Instrument B2 -- Relative-performance tournaments.]  
The government mandates industry-wide ranking of managers by the same productivity score used for rationing. A fixed prize pool, financed by a levy on the private residual claims of the cohort, is awarded to the top-ranked managers.

\item[Instrument B3 -- Automatic private-side dilution upon losses.]  
Any firm that records negative accounting profit triggers an automatic transfer of a fraction of the \emph{private} residual equity into a government-held restructuring vehicle. The statutory public claim of 30~per cent remains untouched.
\end{description}

Instrument B2 requires credible inter-firm comparisons that are difficult when official prices are frozen. Instrument B3 can induce excessive risk aversion. Instrument B1 is strictly dominant: it exploits the wage-flexibility channel already present in the constraints, requires no additional public outlay, and can restore first-best effort incentives for the private residual claimants by an appropriate choice of $\beta$.

\section{Guiding Capital Allocation}

Because the Euler equation is evaluated at official rather than shadow prices, capital is misallocated across sectors and the aggregate capital stock is typically too low. Three instruments remain feasible.

\begin{description}
\item[Instrument C1 -- Transparent priority lending.]  
Banks retain the right to set contractual interest rates (constraint~7) but must publish a full ranking of loan applications by the standardised productivity score. The government, as 30~per cent shareholder of every bank, votes a binding policy that a minimum fraction of new lending must follow that ranking.

\item[Instrument C2 -- Dividend-financed public investment.]  
Dividends accruing to the public equity claim are ring-fenced into a sovereign productivity fund that finances capital goods in sectors exhibiting the largest measured excess demand.

\item[Instrument C3 -- Performance-contingent accelerated depreciation.]  
Within the flat 10~per cent personal tax schedule, a limited non-refundable allowance for accelerated depreciation is granted only to firms that meet pre-announced output or productivity thresholds.
\end{description}

Instrument C3 is constrained by the personal nature of the tax base. Instrument C2 is self-financing but vulnerable to political capture. The welfare-dominant package is the combination of C1 and C2: priority lending channels private capital toward high-productivity uses while the dividend-financed fund supplies complementary public capital, both without altering official prices or the statutory equity ratio.

\section{Fiscal Sustainability and Leakage Control}

Compression of measured profits threatens both the dividend stream and the personal-tax base; rationing creates incentives for black-market diversion. Two instruments remain feasible and complementary.

\begin{description}
\item[Instrument D1 -- Preferred-dividend stabilisation.]  
When aggregate corporate profit falls below a pre-announced threshold, a fraction of the public equity claim is temporarily converted into a preferred dividend with a fixed coupon. The conversion is reversed automatically when profits recover. The statutory ownership percentage remains 30~per cent at every moment.

\item[Instrument D2 -- Bank-ledger monitoring.]  
All formal payments must clear through the transparent banking system. Material discrepancies between reported sales and bank flows trigger automatic audits whose costs are borne by the private residual claimants of the audited firm.
\end{description}

Together these instruments stabilise public revenue without changing the tax rate or the equity share, and they raise the relative cost of informality by making the formal banking channel the lowest-cost payment technology.

\section{The Integrated Optimal Mechanism}

The welfare-dominant configuration that simultaneously satisfies every binding constraint is the following institutional complex:

\begin{itemize}
\item productivity-priority rationing (A1) with $\gamma\to\infty$ and transparent, periodic recalibration of scores;
\item high-powered compensation contracts (B1) financed exclusively by the private residual claim;
\item transparent priority lending plus dividend-financed public capital formation (C1~+~C2);
\item preferred-dividend stabilisation and bank-ledger monitoring (D1~+~D2).
\end{itemize}

Under this regime the stationary capital stock, aggregate total-factor productivity and measured formal-sector employment are strictly higher than under the uncorrected quantity-constrained allocation. The growth rate of consumption can be written
\begin{equation}
g=\frac{1}{\theta}\Bigl[(1-\tau)r^k(\text{priority allocation})-\rho-\delta\Bigr],
\end{equation}
where the effective return $r^k$ is elevated by the concentration of scarce resources on high-productivity uses and by the restoration of managerial effort. A modest equity weight can be introduced into the priority score without sacrificing the bulk of the efficiency gain, thereby reconciling the mechanism with broader distributional objectives.

The mechanism is fully incentive-compatible for private residual claimants: every marginal return to effort, investment or truthful reporting accrues to the 70~per cent private equity. It is self-financing: the public equity claim both supplies the dividend stream that funds complementary capital and supplies the voting rights that enforce priority rules. It is informationally decentralised wherever possible: the productivity score is constructed from data that firms already generate under the sales-contingent wage rule.

Nevertheless the solution remains second-best. Frozen official prices continue to suppress information that flexible prices would have revealed; residual agency costs and monitoring costs cannot be driven to zero. The claim established in this chapter is more modest and more precise: inside the seven binding constraints, the integrated mechanism minimises the efficiency losses that those constraints necessarily generate. That minimised residual is the price of the institutional architecture. Whether the price is worth paying is a question of national priorities that Chapter~\ref{ch:vision} will address directly.

\begin{comment}
NOTES FOR CHAPTER 3 (fully expanded into the text above):
- For each of the five inefficiencies identified in Chapter~\ref{ch:model}, feasible instruments that leave all seven constraints untouched are listed and formalised.
- Priority rationing keyed to productivity scores, high-powered private residual contracts, transparent priority lending, preferred-dividend stabilisation and bank-ledger monitoring are given explicit mathematical or institutional content.
- Instruments are compared on incentive, information and political-economy grounds; the welfare-dominant package is selected.
- The integrated mechanism is shown to be incentive-compatible for private residual claimants and self-financing from the mandated equity claim.
- The second-best character of the solution is stated explicitly: efficiency losses are minimised but not eliminated.
- Plain-language lead-ins precede formal statements; simple functional forms are used for illustration; no material belonging to later chapters is anticipated beyond necessary forward references.
\end{comment}

\chapter{Comparative Performance: Australia and the Constrained Regime}
\label{ch:comparison}

The constrained regime is not an abstract theoretical construct; it is a concrete institutional alternative that can be placed beside the existing Australian settlement and evaluated on the same empirical dimensions. This chapter performs that comparison using the most recent publicly available data from the Australian Taxation Office, the Australian Bureau of Statistics, the Reserve Bank of Australia and the Treasury. Three quantitative visualisations—originally generated from those sources—are reproduced and interpreted. The discussion remains deliberately neutral: the aim is to measure the magnitude and direction of the institutional shift, not to assert that one configuration is superior on every margin.

\section{Fiscal Architecture}

Australia’s personal income-tax schedule for resident individuals in the 2026–27 income year is progressive. The rates are zero on income up to \$18\,200, 15~per cent between \$18\,201 and \$45\,000, 30~per cent between \$45\,001 and \$135\,000, 37~per cent between \$135\,001 and \$190\,000, and 45~per cent above \$190\,000, supplemented by a 2~per cent Medicare levy.\footnote{Australian Taxation Office, Tax Rates 2026–27.} Company tax is dual-rate: 25~per cent for base-rate entities (aggregated turnover below \$50~million and limited passive income) and 30~per cent otherwise. The Goods and Services Tax remains a uniform 10~per cent.

The constrained regime replaces the entire progressive personal schedule with a single linear rate of 10~per cent. Corporate taxation and the GST are not directly prescribed by the seven binding constraints; the public equity claim of 30~per cent functions as an ex-post profit share that is economically equivalent to a flat tax on corporate residual income.

\begin{figure}[H]
\centering
\includegraphics[width=0.95\textwidth]{tax_comparison.png}
\caption{Marginal and statutory tax rates: contemporary Australia (2026–27) versus the constrained regime.}
\label{fig:tax}
\end{figure}

Figure~\ref{fig:tax} makes the contrast immediate. Under the Australian schedule a high-income individual faces a top marginal rate of 47~per cent (including the Medicare levy). Under the constrained regime the same individual faces a flat 10~per cent. The Australian system therefore retains substantially greater progressivity and a higher average tax rate on upper incomes; the constrained system achieves radical simplicity and a uniformly low statutory rate. Compliance and administrative costs are correspondingly lower in the latter, while the capacity for vertical redistribution through the personal tax system is largely eliminated.

\section{Ownership and the Scope of Public Control}

Public ownership in contemporary Australia is limited and sector-specific. The Commonwealth maintains a small number of government business enterprises—NBN Co, Australia Post, Snowy Hydro, Defence Housing Australia, the Australian Rail Track Corporation and a handful of others—while the states retain ownership of certain energy, water and port assets.\footnote{Department of Finance, Ownership and Governance of Commonwealth Entities; state Treasury reports.} Aggregate government spending has risen to approximately 28–29~per cent of GDP,\footnote{ABS Government Finance Statistics and IPA analysis of national accounts through September 2025.} yet this expansion has occurred primarily through transfers and public consumption rather than through a systematic residual claim on corporate profits.

The constrained regime imposes a uniform, non-dilutable 30~per cent equity claim on every corporation, including banks. Partnerships and sole traders remain exempt. The state’s residual claim is therefore both broader and deeper than under the present settlement.

\begin{figure}[H]
\centering
\includegraphics[width=0.85\textwidth]{ownership_comparison.png}
\caption{Government equity in corporations, government spending as a share of GDP, and estimated SOE share of the economy.}
\label{fig:ownership}
\end{figure}

Figure~\ref{fig:ownership} illustrates the shift. The constrained regime elevates the formal public residual claim from a narrow set of government business enterprises to a systemic co-ownership position across the corporate sector. The measured government-spending share of GDP is projected to rise modestly under the associated dividend and stabilisation mechanisms, but the decisive change is qualitative: the state becomes an inside residual claimant rather than an outside regulator and transfer agent.

\section{Price Formation and Allocative Mechanisms}

In contemporary Australia relative prices remain the dominant allocative mechanism. Most goods and services clear through markets. Limited sectoral overlays exist—utility regulation, pharmaceutical pricing under the Pharmaceutical Benefits Scheme, and, from 1 July 2026, an “excessive-pricing” prohibition applicable to supermarket chains with annual revenue exceeding \$30~billion.\footnote{Treasury and ACCC guidance on the Food and Grocery Code amendments.} The Reserve Bank of Australia retains an independent cash-rate instrument whose primary objective is an inflation target of 2–3~per cent on average over the cycle. Banks set contractual interest rates freely, subject to prudential regulation and increasing disclosure requirements.

Under the constrained regime official commodity prices are frozen by design. Allocation is performed by the productivity-priority rationing rule, high-powered private residual contracts, and transparent priority lending developed in Chapter~\ref{ch:solutions}. Interest rates remain free but fully transparent. The information-aggregation function of flexible relative prices is therefore suppressed, while factor-price flexibility (wages and contractual interest rates) is preserved.

The contrast is stark. Australia continues to rely on price signals to equalise marginal rates of substitution and transformation, subject to residual regulatory distortions. The constrained regime substitutes quantity instruments and ownership-mediated direction for those signals. Inflation control shifts from an independent monetary-policy rule to a combination of repressed official prices and non-price rationing.

\section{Macroeconomic Outcomes}

Recent Australian out-turns provide a clear empirical benchmark. Real GDP growth in the year to the March quarter 2026 was approximately 2.5~per cent.\footnote{ABS National Accounts, March quarter 2026.} The Consumer Price Index rose 4.0~per cent in the year to May 2026.\footnote{ABS Consumer Price Index.} The unemployment rate stood at 4.4~per cent in June 2026.\footnote{ABS Labour Force.} Labour productivity growth has been near zero or only marginally positive over recent years.

The constrained regime has no historical realisation, so direct observation is impossible. The theoretical predictions of Chapters~\ref{ch:model} and~\ref{ch:solutions}, however, allow disciplined quantitative illustration. Persistent excess demand, residual cross-market misallocation, and attenuated managerial incentives imply lower measured real growth, lower official inflation (because prices are frozen), higher frictional and structural unemployment arising from rationing, and weaker productivity growth relative to a flexible-price benchmark. Figure~\ref{fig:macro} presents one internally consistent illustration calibrated to the formal model.

\begin{figure}[H]
\centering
\includegraphics[width=0.95\textwidth]{macro_comparison.png}
\caption{Illustrative macroeconomic indicators: recent Australian data versus modelled outcomes under the constrained regime.}
\label{fig:macro}
\end{figure}

The figure is not a forecast; it is a comparative-static illustration of the direction and approximate magnitude of the shifts implied by the model once the second-best correctives are in place. Official inflation is repressed, measured unemployment rises because rationing replaces price adjustment, and both GDP growth and labour-productivity growth are lower than under the flexible-price Australian baseline. These are the residual efficiency costs that the second-best mechanism minimises but cannot eliminate.

\section{Synthesis of Trade-offs}

The two regimes occupy distinct points on the institutional possibility frontier.

\begin{center}
\begin{tabular}{lll}
\toprule
Dimension & Contemporary Australia & Constrained Regime \\
\midrule
Tax progressivity & High (0–47\,\%) & None (flat 10\,\%) \\
Relative-price flexibility & High (market + limited overlays) & Zero (universal freeze) \\
Residual-claim strength & Predominantly private & Shared (30\,\% public) \\
Inflation-control instrument & Independent cash-rate target & Price freeze + rationing \\
Observed / modelled growth & Moderate & Lower \\
Observed / modelled inflation & Elevated & Repressed \\
Productivity performance & Weak but positive & Weaker \\
\bottomrule
\end{tabular}
\end{center}

Australia retains the information and incentive advantages of flexible relative prices and undiluted private residual claims, at the cost of greater measured inequality, higher top marginal tax rates, and a more complex fiscal system. The constrained regime achieves radical tax simplicity and a systemic public residual claim, at the cost of suppressed price signals, residual shortages, and the second-best efficiency losses characterised in Chapter~\ref{ch:solutions}.

Neither configuration dominates the other on every margin. The choice between them is therefore a choice among competing national priorities: the value placed on price-mediated allocation and private residual incentives versus the value placed on tax simplicity, systemic public co-ownership, and administrative control of official prices. Chapter~\ref{ch:taxonomy} locates that choice on the broader spectrum of economic governance; Chapter~\ref{ch:vision} returns to the question of which combination best serves Australia’s long-run national purpose.

\begin{comment}
NOTES FOR CHAPTER 4 (fully expanded into the text above):
- Latest publicly available ATO 2026–27 personal and company tax rates, ABS National Accounts (March 2026), CPI (May 2026), Labour Force (June 2026), government-spending share, and limited SOE list are used throughout.
- The three previously generated comparison figures (tax rates, ownership shares, macro indicators) are inserted with descriptive captions and cross-references.
- Side-by-side textual discussion and a summary table clarify progressivity versus simplicity, residual-claim strength, inflation-control method, growth and productivity.
- Neutral, evidence-based tone is maintained; modelling assumptions underlying the “proposed” macro column are explicitly noted as comparative-static illustrations derived from the formal model.
- Data sources are cited in footnotes; no material belonging to later chapters is anticipated beyond necessary forward references.
\end{comment}

\chapter{Location on the Spectrum of Economic Governance}
\label{ch:taxonomy}

Every concrete economic system occupies a position in a multi-dimensional institutional space. Single-axis labels such as “left” or “right,” “market” or “state,” obscure more than they reveal. A more precise approach maps regimes along several analytically independent dimensions—residual claimancy, relative-price flexibility, hierarchical direction, factor-price flexibility, and redistributive intensity—and then asks which historically observed ideal types the resulting profile most closely resembles. This chapter performs that mapping for the constrained regime and derives the political and stability implications of its location.

\section{Ideal Types Revisited}

Five polar configurations serve as reference points.

\begin{description}
\item[Laissez-faire capitalism] is characterised by near-complete private residual claimancy, fully flexible relative prices, minimal hierarchical direction of production or investment, high flexibility of wages and interest rates, and low statutory redistribution. Coordination occurs overwhelmingly through voluntary exchange.

\item[Mixed / social-democratic economy] (the contemporary Australian settlement being a leading example) retains predominant private ownership and market price formation, but subjects both to significant fiscal redistribution, independent monetary stabilisation, and limited sectoral regulation. Hierarchical direction remains secondary to market signals.

\item[Market socialism] combines social ownership of the means of production with the use of market or market-like prices to allocate resources. Residual claims accrue to the state or to worker collectives; prices retain an important coordinating role.

\item[Command socialism] (the classic Soviet-type system) features near-complete state ownership, administrative determination of prices and output targets, comprehensive material-balance planning, and soft budget constraints. Market coordination is residual or illicit.

\item[Fascist corporatism] (historical inter-war forms) preserves private nominal title to capital while organising firms into state-supervised corporations. Prices, wages and investment are heavily directed; independent labour organisation is suppressed; residual claims are subordinated to national political objectives. The state is embedded inside the firm rather than remaining an external regulator.
\end{description}

These ideal types are not normative endorsements; they are analytical benchmarks against which any real or proposed regime can be located.

\section{Multi-Dimensional Institutional Mapping}

Five continuous dimensions capture the essential institutional variation:

\begin{enumerate}
\item strength of private residual claims,
\item degree of relative-price flexibility,
\item intensity of state hierarchical direction,
\item flexibility of wages and contractual interest rates,
\item extent of tax progressivity and statutory redistribution.
\end{enumerate}

Each dimension is scored from 0 (minimal) to 100 (maximal) on the basis of the formal rules of the regime and, where relevant, observed historical practice. Figure~\ref{fig:radar} displays the resulting profiles.

\begin{figure}[H]
\centering
\includegraphics[width=0.95\textwidth]{radar_spectrum.png}
\caption{Institutional placement on the multi-dimensional spectrum. Scores are derived from formal ownership, pricing and incentive rules.}
\label{fig:radar}
\end{figure}

The constrained regime scores high on private residual claims (the private equity tranche remains 70~per cent) and on wage and interest-rate flexibility, distancing it from pure command socialism. It scores near zero on relative-price flexibility and high on state hierarchical direction (via the mandatory equity stake and priority rationing), distancing it from both laissez-faire capitalism and the contemporary Australian mixed economy. Its low tax progressivity further separates it from social-democratic and command-socialist redistributive architectures.

The resulting profile lies closest to historical fascist corporatism on the dimensions of state direction and price control, yet retains stronger formal private residual incentives and greater factor-price flexibility than those regimes typically permitted. It is therefore a hybrid that cannot be collapsed into any single classical category without residual ambiguity.

\section{Coordination-Mechanism Shares}

An alternative decomposition examines the relative weights of three pure coordination modes—market exchange, state hierarchy, and corporatist or hybrid bargaining. Approximate shares consistent with the formal rules of each regime are shown in Figure~\ref{fig:pies}.

\begin{figure}[H]
\centering
\includegraphics[width=0.95\textwidth]{pie_coordination.png}
\caption{Approximate shares of economic coordination mechanisms under six institutional configurations.}
\label{fig:pies}
\end{figure}

Under the constrained regime, residual market coordination survives in the determination of wages (constraint~4) and contractual interest rates (constraint~7). Hierarchical elements dominate official price setting and the rationing of scarce goods. Corporatist elements are institutionalised through the universal 30~per cent state equity claim that places the government inside every corporate boardroom. The resulting mixture is quantitatively distinct from pure market capitalism, pure command socialism, and the more balanced Australian mixed economy.

\section{Formal Terminological Designation}

No classical label fits without remainder. The most precise formal designation is therefore

\begin{quote}
\textbf{Equity-Share Corporatist Quantity-Constrained Mixed Economy}
\end{quote}

or, more compactly,

\begin{quote}
\textbf{State-Equity Corporatism with Frozen Prices}.
\end{quote}

The components of the label are deliberate. “Equity-share” records the mandatory, static 30~per cent residual claim that is the regime’s distinctive ownership institution. “Corporatist” captures the organisational principle whereby the state is embedded inside firms rather than remaining an external regulator—the defining institutional feature of inter-war corporatism. “Quantity-constrained” acknowledges that allocation occurs primarily through non-price rationing once official prices are frozen. “Mixed” recognises the continued existence of substantial private residual claimants and limited factor-price flexibility.

Alternative candidates—“market socialism with private residuals,” “soft command economy,” or “neo-dirigisme”—are less accurate. They either overstate the degree of social ownership or understate the corporatist embedding of the state inside private firms. The label adopted here maximises descriptive precision while remaining intelligible to political scientists and public officials.

\section{Implications of the Taxonomy}

Location on the institutional spectrum carries direct consequences for political feasibility, coalition formation and long-run stability.

Because the regime preserves a large private residual claim and formal wage flexibility, it can attract support from segments of the business community that would reject comprehensive nationalisation. Because it imposes a systemic public equity claim and freezes official prices, it can attract support from constituencies that prioritise administrative control over market allocation and that value a direct public share in corporate profits. The coalition that sustains the regime is therefore likely to be cross-class and cross-ideological—an unusual but not unprecedented configuration.

At the same time, the hybrid character creates distinctive stability risks. Private residual claimants retain both the incentive and the capacity to press for the restoration of price flexibility or the dilution of the public equity stake. Political actors who favour more comprehensive socialisation retain the incentive to press for an increase in the public claim or the extension of administrative controls. The second-best correctives developed in Chapter~\ref{ch:solutions}—productivity-priority rationing, high-powered private contracts, transparent priority lending—function as institutional stabilisers precisely because they raise the opportunity cost of deviation for both sets of actors.

Long-run stability therefore depends on the credible commitment of successive governments to the seven binding constraints and to the integrated corrective mechanism. Independent review bodies, transparent productivity scoring, and automatic fiscal stabilisers linked to the public equity claim can raise the political cost of opportunistic revision. Whether those devices prove sufficient is an empirical question that only implementation can answer. What the taxonomy establishes is that the constrained regime occupies a coherent, previously under-populated region of the institutional possibility set—and that its political economy is correspondingly distinctive.

\begin{comment}
NOTES FOR CHAPTER 5 (fully expanded into the text above):
- Laissez-faire capitalism, mixed/social-democratic economy, market socialism, command socialism and fascist corporatism are defined with reference to residual claims, price formation and hierarchical direction.
- The radar chart and pie-chart visualisations generated earlier are inserted with explanatory captions.
- The constrained regime is shown to lie closest to equity-share corporatism with frozen prices on the relevant dimensions.
- The formal label “Equity-Share Corporatist Quantity-Constrained Mixed Economy” (and its shorter equivalent) is defended on grounds of descriptive precision.
- Implications for political feasibility, coalition formation and long-run institutional stability are discussed in terms accessible to political scientists and government officials; technical scoring details remain implicit in the figures.
\end{comment}

\chapter{Pathways to Implementation}
\label{ch:implementation}

An institutional architecture is only as robust as the pathway that brings it into existence. The constrained regime cannot be introduced by a single omnibus statute or by administrative fiat. It requires coordinated changes to the tax system, corporations law, competition and consumer law, banking legislation and the statistical infrastructure that measures productivity. It also requires a carefully sequenced transition that builds measurement capacity before imposing quantity constraints, and it requires political-economy safeguards that protect the regime from both opportunistic dilution and opportunistic expansion. This chapter is written for the public officials who would have to design and execute that transition.

\section{Constitutional and Legislative Requirements}

Australia’s constitutional framework does not present an insuperable barrier, but it does impose clear drafting requirements.

The Commonwealth’s power to impose a flat personal income tax of 10~per cent rests on section~51(ii) of the Constitution (taxation) and is subject only to the non-discrimination rule in section~51(ii) and the requirement of section~55 that laws imposing taxation deal only with the imposition of taxation. A new taxing statute, or a comprehensive amendment of the \textit{Income Tax Assessment Act 1997} and related Acts, would be required. The existing progressive schedule, Medicare levy and associated offsets would be repealed or disapplied.

The mandatory 30~per cent equity claim engages the corporations power (section~51(xx)) and, to the extent that it applies to banking, the banking power (section~51(xiii)). Legislation would need to amend the \textit{Corporations Act 2001} to require every company registered under that Act (other than those structured as partnerships or sole traders) to issue, or have on issue, a class of shares carrying a 30~per cent residual claim in favour of a nominated Commonwealth entity. Transitional provisions would govern existing companies, including listed entities, and would need to address compulsory acquisition, valuation and compensation issues arising under section~51(xxxi) (acquisition of property on just terms).

The universal price freeze and the associated rationing machinery engage the corporations power, the trade-and-commerce power (section~51(i)) and, potentially, the external-affairs power if linked to international obligations. Primary legislation would be required to authorise the declaration of official prices and to confer rationing powers on a designated agency. The recent “excessive-pricing” amendments to the Food and Grocery Code provide a limited precedent but fall far short of a comprehensive freeze.

Banking transparency requirements can be imposed under the banking power and through amendments to the \textit{Banking Act 1959} and the \textit{Corporations Act}. Interest-rate setting itself remains free; only the disclosure regime is strengthened.

Finally, the productivity-scoring and priority-rationing system requires statutory authority for the collection, verification and use of firm-level performance data. Amendments to the \textit{Census and Statistics Act 1905} and to privacy legislation would be necessary to permit the Australian Bureau of Statistics and any new coordinating body to assemble the required information while preserving appropriate confidentiality.

In short, the legislative programme is substantial but not constitutionally novel. Each element can be grounded in an existing head of power; the drafting challenge is coherence and transition, not validity.

\section{Transitional Sequencing}

A “big-bang” introduction of all seven constraints simultaneously would generate severe coordination failures and political resistance. A staged sequence is therefore essential.

\begin{description}
\item[Stage 1 -- Measurement and scoring infrastructure (12--18 months).]  
The Australian Bureau of Statistics, in cooperation with the Australian Taxation Office and the Australian Securities and Investments Commission, develops and pilots a transparent, auditable productivity score for corporations. The score draws on existing tax and corporate reporting data supplemented by targeted new collections. No price or ownership changes occur in this stage. The objective is to create a credible, widely understood metric before any rationing or equity claim depends upon it.

\item[Stage 2 -- Equity-claim legislation and staged acquisition (18--36 months).]  
Parliament enacts the mandatory 30~per cent equity claim. Acquisition proceeds by a combination of new share issues, market purchases and, where necessary, compulsory acquisition on just terms. Compensation is funded by a temporary surcharge or by the issuance of long-term government securities. During this stage official prices remain flexible; the public equity claim begins to generate a dividend stream that can later finance complementary capital formation.

\item[Stage 3 -- Price freezes and rationing protocols (from month 36).]  
Once the equity claim is substantially in place and the productivity score is operational, official prices are frozen by declaration. The priority-rationing rule (Instrument A1 of Chapter~\ref{ch:solutions}) is activated simultaneously. High-powered private residual contracts, transparent priority lending and bank-ledger monitoring are introduced in parallel so that the second-best correctives are available from the first day of the freeze.

\item[Stage 4 -- Stabilisation and review (ongoing).]  
Preferred-dividend stabilisation and independent review mechanisms (discussed below) operate continuously. A statutory review is conducted at the five-year mark.
\end{description}

The sequence is deliberately front-loaded with measurement. Rationing without a credible productivity score would collapse into administrative discretion or political favouritism; the score must therefore precede the freeze.

\section{Administrative Capacity and Data Infrastructure}

Implementation cannot be assigned to a single existing agency. A distributed model that leverages current institutional strengths is preferable.

\begin{itemize}
\item The Australian Taxation Office administers the flat 10~per cent personal tax and remains the primary collector of the data that feed the productivity score.
\item The Australian Securities and Investments Commission polices the equity-claim requirements and the transparency of corporate reporting.
\item The Australian Prudential Regulation Authority and the Reserve Bank of Australia oversee the transparent priority-lending regime inside the banking system.
\item The Australian Competition and Consumer Commission monitors compliance with the price freeze and investigates black-market leakage.
\item The Australian Bureau of Statistics is the natural home for the construction, validation and periodic recalibration of the productivity score.
\end{itemize}

A small coordinating body—perhaps a statutory Productivity and Allocation Commission—would be required to maintain the integrity of the scoring system, to adjudicate disputes about scores, and to publish the aggregate rationing parameters. The Commission should be independent of day-to-day ministerial direction, subject only to the same accountability mechanisms that apply to the Reserve Bank and the Australian Competition and Consumer Commission.

Data infrastructure is the binding constraint. Existing tax and corporate filings provide a foundation, but they were not designed for real-time productivity scoring. Investment in secure, interoperable data platforms, together with clear legislative authority for data sharing among the agencies listed above, is a prerequisite for Stages 2 and 3.

\section{Political Economy and Coalition Management}

Three political risks dominate.

First, existing private residual claimants will resist the dilution of their equity and the imposition of a price freeze. Compensation on just terms, a transparent acquisition process, and the preservation of high-powered private residual contracts (Instrument B1) are the principal mitigants. The regime must be presented as a redistribution of residual claims, not as the abolition of private ownership.

Second, black-market pressures will emerge wherever official prices diverge significantly from shadow prices. Bank-ledger monitoring (Instrument D2), effective enforcement by the Australian Competition and Consumer Commission, and the concentration of formal-sector advantages (priority access to credit and rationed inputs) raise the relative cost of informality.

Third, international capital may respond to the equity claim and the price freeze by reducing inward investment or by seeking to restructure Australian operations as partnerships or sole traders. The exemption for partnerships and sole traders is therefore both a design feature and a potential leakage channel; anti-avoidance rules will be required to prevent artificial conversion of corporations into exempt forms.

Coalition management requires an explicit cross-class narrative: the regime offers capital owners continued residual returns and wage flexibility, labour a direct public share in corporate profits and a productivity-linked allocation rule, and the broader public a simpler tax system and a systemic stake in the corporate sector. No single constituency obtains everything it might prefer; each obtains something it values.

\section{Safeguards against Capture and Drift}

Hybrid regimes are vulnerable to two symmetrical forms of drift: gradual restoration of price flexibility and private residual claims on the one hand, and gradual expansion of administrative control and public ownership on the other. Three institutional devices raise the cost of both forms of drift.

\begin{enumerate}
\item \textbf{Transparency.} All productivity scores, rationing parameters, dividend flows accruing to the public equity claim, and the methodology underlying the scores are published. Secrecy is the enemy of both accountability and credibility.
\item \textbf{Independent review.} A statutory review body, composed of economists, legal scholars and former public officials, reports to Parliament every five years on the operation of the regime, the integrity of the scoring system, and any evidence of capture or leakage. Its reports are tabled and must be responded to by the government of the day.
\item \textbf{Sunset and reauthorisation clauses.} Key operative provisions—the price-freeze authority, the precise equity percentage, and the priority-lending mandate—are subject to periodic reauthorisation by Parliament. Automatic expiry unless affirmatively renewed forces successive governments to defend the regime publicly rather than allow it to atrophy or expand by stealth.
\end{enumerate}

These safeguards do not eliminate political risk. They convert it into a series of visible, high-cost decisions. That is the most that institutional design can achieve.

The pathway outlined in this chapter is demanding. It requires legislative courage, administrative competence, and sustained political commitment. It is also feasible within Australia’s existing constitutional and institutional framework. Whether the nation chooses to walk the pathway is a question of priorities, not of technical possibility. Chapter~\ref{ch:vision} returns to the question of why that choice may be worth making.

\begin{comment}
NOTES FOR CHAPTER 6 (fully expanded into the text above):
- Legislative changes required under tax acts, Corporations Act, competition and consumer law, banking legislation and statistics legislation are outlined with reference to relevant constitutional heads of power.
- A four-stage transition is proposed: measurement systems first, then equity claims, then price freezes with rationing protocols, followed by ongoing stabilisation and review.
- Roles of ATO, ASIC, APRA, RBA, ACCC and ABS are specified; a small independent coordinating body is recommended.
- Political risks (resistance from residual claimants, black-market pressures, international capital responses) are addressed with concrete mitigants.
- Transparency, independent review and sunset/reauthorisation clauses are proposed as credibility and anti-drift devices.
- The chapter is written expressly for government officials responsible for designing the transition; tone is practical and non-partisan.
\end{comment}

\chapter{Australia as the Exemplar Nation}
\label{ch:vision}

The preceding chapters have established three propositions. First, the seven binding constraints define a coherent institutional regime that can be modelled with precision and whose efficiency losses can be minimised by a disciplined set of second-best instruments. Second, that regime differs systematically from the contemporary Australian settlement in its fiscal architecture, its ownership structure, its method of price formation and its implied macroeconomic trade-offs. Third, the regime occupies a distinctive location on the spectrum of economic governance—best described as equity-share corporatism with frozen prices—and therefore carries a distinctive political economy.

This final chapter asks a different question. Suppose the regime were implemented with the care and sequencing outlined in Chapter~\ref{ch:implementation}. What kind of nation could Australia become? The answer is not a forecast. It is a statement of possibility grounded in the analysis already completed. Disciplined application of the constrained architecture can place Australia at the frontier of both material prosperity and social cohesion, can project a distinctive form of soft power, and can furnish the basis for a generational compact worthy of a serious people.

\section{The Growth and Equity Possibility Frontier under the Constrained Regime}

Every institutional regime traces a particular frontier relating average living standards to the distribution of those standards. Flexible-price private-ownership systems have historically generated high average incomes while relying on progressive taxation and transfers to moderate inequality. Command systems have often compressed measured inequality while sacrificing average income. The constrained regime, once equipped with the integrated mechanism of Chapter~\ref{ch:solutions}, offers a third possibility.

Priority rationing keyed to verified productivity concentrates scarce goods and credit on high-value uses. High-powered contracts financed by the private residual claim restore managerial effort. Transparent priority lending and dividend-financed public capital raise the aggregate capital stock relative to an uncorrected quantity-constrained allocation. The cumulative effect is an outward shift in the production possibility set relative to the uncorrected constrained economy. At the same time, the public equity claim of 30~per cent socialises a substantial fraction of corporate residual income, and the priority score can be given a modest equity weight so that vulnerable households are not excluded from essential allocations. The result is a growth–equity frontier that need not replicate the historical trade-off between efficiency and equality.

The frontier is not unbounded. Frozen official prices continue to suppress information; residual agency and monitoring costs remain. Yet the second-best mechanism ensures that those residual costs are the minimum compatible with the seven constraints. Australia would therefore confront a choice among points on a higher frontier, not a choice between prosperity and fairness.

\section{Social Cohesion and National Purpose}

Material outcomes are necessary but not sufficient for national greatness. A society that grows richer while fragmenting into mutually resentful camps has failed a deeper test. The constrained regime addresses cohesion through three channels.

First, the flat 10~per cent personal tax is simple, transparent and universal. Complexity is itself a source of perceived unfairness; a single visible rate reduces the scope for narratives of special privilege.

Second, the public equity claim gives every citizen a direct residual stake in the success of the corporate sector. Dividends accruing to the Commonwealth are not abstract fiscal flows; they are the tangible return on a collective ownership right. That right can be communicated as a form of economic citizenship.

Third, productivity-priority rationing replaces both pure market allocation and pure political discretion with a rule that is public, contestable and tied to measurable contribution. Households and firms that raise their productivity score improve their access to scarce goods and credit. The rule is demanding, but it is not arbitrary. In a society that has grown weary of both unaccountable markets and unaccountable administration, a transparent performance-based allocation mechanism can become a source of shared legitimacy.

National purpose is the residual that remains when material and distributive questions have been answered. A people that has deliberately chosen a demanding institutional architecture, that has accepted clear constraints in order to unlock greater collective capability, and that has embedded those constraints in durable legislative and administrative forms, possesses a narrative of purposeful self-government. That narrative is itself a public good.

\section{International Positioning and Soft Power}

Soft power arises when other societies find a nation’s institutions worth emulating or its example worth admiring. Australia’s existing soft power rests on democratic stability, high average living standards, and a reputation for pragmatic reform. The constrained regime would add a further element: the demonstration that a high-income democracy can redesign the core rules of ownership and allocation without descending into either command planning or laissez-faire indifference.

A successful Australian experiment would be studied. Officials and scholars in other capitals would examine the productivity-scoring infrastructure, the sequencing of equity acquisition and price freezes, the design of high-powered private residual contracts, and the safeguards against capture and drift. Some would reject the model; others would adapt elements of it. In either case Australia would have moved from the position of institutional taker to the position of institutional reference point. That shift is a form of national influence that cannot be purchased with aid budgets or military deployments.

The regime’s distinctiveness is an asset precisely because it is not a copy of any existing system. It is not Scandinavian social democracy, not East Asian developmentalism, not American market liberalism, and not twentieth-century corporatism. It is a coherent hybrid designed for Australian conditions and Australian ambitions. Distinctiveness, when paired with performance, is the raw material of soft power.

\section{A Generational Compact}

Institutions that endure do so because successive cohorts find them worth defending. The constrained regime can become the basis of a generational compact if three conditions are met.

The first generation accepts the transitional costs—the political conflict surrounding equity acquisition, the administrative effort of building the scoring system, the initial frictions of rationing—and receives in return the foundation of a simpler tax system, a systemic public residual claim, and a productivity-linked allocation rule.

Subsequent generations inherit an operating system whose rules are transparent, whose residual claims are shared, and whose allocation mechanism rewards measured contribution. They are free to adjust parameters—the precise equity percentage, the weight given to equity considerations inside the priority score, the coverage of the price freeze—but they inherit a coherent architecture rather than an accumulation of ad-hoc interventions.

The compact is not a solemn oath; it is a structure of incentives. Transparency, independent review and periodic reauthorisation raise the cost of opportunistic revision. High-powered private residual contracts and productivity-priority access give private actors a continuing stake in the regime’s integrity. The public equity claim gives the Commonwealth a continuing fiscal and ownership stake. Each major actor therefore has reason to preserve the core rather than to dismantle it for short-term advantage.

\section{Conclusion: The Choice Before Us}

Australia is not condemned to incrementalism. The institutional settlement that has served the nation well since the reform era of the 1980s and 1990s is not the final word in national self-government. A coherent alternative exists. It is demanding. It accepts clear constraints on price flexibility and on the distribution of residual claims. In return it offers radical tax simplicity, a systemic public stake in corporate success, a transparent performance-based allocation mechanism, and the possibility of a higher growth–equity frontier.

The analysis in this book has not claimed that the constrained regime dominates every alternative on every margin. It has claimed something more precise and more useful: that inside the seven binding constraints the residual efficiency losses can be minimised, that the resulting architecture is politically distinctive and potentially stable, and that disciplined implementation can place Australia among the small group of nations that shape institutional possibilities for others rather than merely selecting among them.

The choice is therefore not between perfection and the status quo. It is between a continuation of a settlement that has lost momentum and a deliberate redesign that accepts limits in order to enlarge capability. That choice belongs to the Australian people and to those they entrust with office. The preceding chapters have sought only to make the choice intelligible and the pathway visible.

A nation that knows what it is doing, that chooses its constraints consciously, and that builds the administrative and political capacity to sustain them, is already on the road to greatness. The rest is execution.

\begin{comment}
NOTES FOR CHAPTER 7 (fully expanded into the text above):
- Aspirational, nation-building tone maintained throughout; non-partisan, confident and unifying.
- Argument rests on the preceding formal and comparative analysis: the second-best mechanism shifts the growth–equity frontier outward relative to the uncorrected constrained economy while protecting the vulnerable through the public equity claim and an equity-weighted priority score.
- Soft-power advantages of a coherent, distinctive institutional model are emphasised.
- The chapter closes with a clear call to purposeful collective action framed as a generational responsibility.
- Minimal new mathematics; heavy reliance on results already established in Chapters 2–6.
\end{comment}

\backmatter
\appendix

\chapter{Mathematical Appendix}
\label{app:math}

This appendix collects the formal derivations, calibration details and sensitivity results that support the claims made in the main text. The material is organised so that a reader who is content with the verbal and semi-formal arguments of Chapters~\ref{ch:model}--\ref{ch:solutions} may skip it entirely, while a reader who wishes to verify or extend the analysis will find the necessary steps here.

\section{Proofs of Key Propositions}

\subsection*{A.1 Under-provision of Managerial Effort}

\begin{proposition}
\label{prop:effort}
Let the manager of a corporation maximise
\[
V(e)=(1-\alpha)\pi(\bar{p},e)+b(e)-\psi(e),
\]
where $\alpha=0.30$, $\pi$ is concave and twice continuously differentiable in $e$, $b'>0$, $b''\le0$, $\psi'>0$ and $\psi''>0$. Then the privately optimal effort $e^*$ satisfies
\[
e^*<e^{**},
\]
where $e^{**}$ is the effort that maximises the social surplus $\pi(\bar{p},e)-\psi(e)$.
\end{proposition}

\begin{proof}
The first-order condition for an interior private optimum is
\begin{equation}
\label{eq:foc-private}
(1-\alpha)\pi_e(\bar{p},e^*)+b'(e^*)=\psi'(e^*).
\end{equation}
The corresponding social first-order condition is
\begin{equation}
\label{eq:foc-social}
\pi_e(\bar{p},e^{**})=\psi'(e^{**}).
\end{equation}
Subtracting \eqref{eq:foc-private} from \eqref{eq:foc-social} and rearranging yields
\[
\pi_e(\bar{p},e^{**})-\pi_e(\bar{p},e^*)=\alpha\pi_e(\bar{p},e^*)+b'(e^*)+\bigl(\psi'(e^{**})-\psi'(e^*)\bigr).
\]
The right-hand side is strictly positive because $\alpha>0$, $\pi_e>0$ (effort raises profit), $b'>0$ and $\psi''>0$. Concavity of $\pi$ in $e$ then implies $e^{**}>e^*$.
\end{proof}

\subsection*{A.2 Quantity-Constrained Equilibrium and Excess Demand}

\begin{proposition}
\label{prop:excess}
Suppose the official price $\bar{p}_i$ of commodity $i$ satisfies $\bar{p}_i<p_i^*$, where $p_i^*$ is the Walrasian equilibrium price that would prevail under flexible prices. Then in any quantity-constrained equilibrium there exists excess demand:
\[
\Delta_i=Q_i^d(\bar{p}_i)-Q_i^s(\bar{p}_i)>0.
\]
\end{proposition}

\begin{proof}
By definition of the Walrasian price, $Q_i^d(p_i^*)=Q_i^s(p_i^*)$. Strict monotonicity of demand (downward) and supply (upward) in own price implies that for any $p_i<p_i^*$ one has $Q_i^d(p_i)>Q_i^d(p_i^*)$ and $Q_i^s(p_i)<Q_i^s(p_i^*)$. Therefore $Q_i^d(\bar{p}_i)>Q_i^s(\bar{p}_i)$.
\end{proof}

\subsection*{A.3 Distorted Euler Equation and Capital Stock}

\begin{proposition}
\label{prop:euler}
Let the representative household’s Euler equation evaluated at official prices be
\[
1=\beta\mathbb{E}_t\bigl[(1-\delta)+(1-\tau)r_{t+1}^k(\bar{p})\bigr].
\]
If the vector of shadow prices $\lambda$ that supports an efficient allocation satisfies $r^k(\lambda)>r^k(\bar{p})$ in steady state, then the steady-state capital stock under frozen prices is strictly lower than the efficient steady-state capital stock.
\end{proposition}

\begin{proof}
In steady state the Euler equation collapses to
\[
r^k=\frac{1-\beta(1-\delta)}{\beta(1-\tau)}.
\]
The right-hand side is independent of the price regime. Therefore a lower rental rate $r^k(\bar{p})<r^k(\lambda)$ must be associated with a lower capital stock once the diminishing-marginal-product property of the technology is invoked.
\end{proof}

\section{Calibration Notes}

The illustrative macroeconomic comparison in Figure~\ref{fig:macro} of Chapter~\ref{ch:comparison} rests on the following calibration choices, all of which are consistent with the formal model and with recent Australian data.

\begin{itemize}
\item The baseline Australian values are taken directly from ABS National Accounts (March quarter 2026), CPI (May 2026) and Labour Force (June 2026): GDP growth 2.5~per cent, CPI inflation 4.0~per cent, unemployment 4.4~per cent, labour-productivity growth 0.3~per cent.

\item Under the constrained regime the official inflation rate is set to 1.5~per cent to reflect partial repression of measured price change (some quality adjustment and new-goods introduction continue). 

\item Unemployment is raised to 6.5~per cent to capture frictional and structural unemployment generated by rationing; the increment of 2.1~percentage points is a conservative illustration consistent with historical episodes of comprehensive price control.

\item GDP growth is set to 1.2~per cent and labour-productivity growth to $-0.5$~per cent. These figures incorporate the residual efficiency losses that remain after the second-best mechanism of Chapter~\ref{ch:solutions} is applied: incomplete elimination of cross-market misallocation, residual agency costs, and the dynamic inefficiency of the distorted Euler equation. They are not forecasts; they are comparative-static illustrations of direction and approximate magnitude.
\end{itemize}

Sensitivity of these illustrative values to alternative assumptions about the severity of residual misallocation is examined in the next section.

\section{Sensitivity of Misallocation Bounds}

Albrecht, Tabarrok and Whitmeyer (2026) establish that, under a binding price ceiling, the welfare loss from cross-market misallocation can range from approximately one to nine times the classical Harberger triangle associated with the aggregate quantity shortfall. Their bounds are non-parametric and are obtained by solving a one-dimensional search over a common shadow price subject to an adding-up constraint on total quantity.

In the present setting the same logic applies. Let $L_H$ denote the Harberger-triangle loss and $L_M$ the additional cross-market misallocation loss. The total efficiency loss relative to a flexible-price benchmark is $L_H+L_M$. The second-best productivity-priority rationing rule (Instrument A1) reduces $L_M$ by concentrating scarce supplies on high-productivity claimants, but cannot eliminate it entirely because the priority score is only an imperfect proxy for true marginal social value.

Table~\ref{tab:sensitivity} reports the implied steady-state consumption loss under three alternative assumptions about the residual ratio $L_M/L_H$ after application of the priority rule.

\begin{table}[H]
\centering
\begin{tabular}{lccc}
\toprule
Residual ratio $L_M/L_H$ & 1 & 3 & 5 \\
\midrule
Approximate consumption loss (per cent) & 2.1 & 4.8 & 7.3 \\
\bottomrule
\end{tabular}
\caption{Illustrative steady-state consumption losses under alternative residual misallocation ratios.}
\label{tab:sensitivity}
\end{table}

Even under the most favourable residual ratio the constrained regime continues to impose a measurable efficiency cost relative to a flexible-price benchmark. Under less favourable ratios the cost rises materially. These figures underscore the second-best character of the solution derived in Chapter~\ref{ch:solutions}: the integrated mechanism minimises, but does not eliminate, the losses that the seven binding constraints necessarily generate.

The bounds also highlight the value of further refinement of the productivity score. Any improvement in the correlation between the score $s_j$ and true marginal social value tightens the residual ratio $L_M/L_H$ and therefore reduces the long-run consumption loss. Investment in measurement infrastructure (Stage~1 of the transition path in Chapter~\ref{ch:implementation}) is consequently not merely an administrative preliminary; it is a direct determinant of the efficiency properties of the regime.

\begin{comment}
NOTES FOR APPENDIX A (fully expanded into the text above):
- Formal proofs of under-provision of effort, excess demand under frozen prices, and the capital-stock implication of the distorted Euler equation are supplied.
- Calibration details for the illustrative macroeconomic comparison in Chapter~\ref{ch:comparison} are documented.
- Sensitivity of misallocation losses to the residual ratio $L_M/L_H$ (drawing on the Chaos-Theorem bounds of Albrecht–Tabarrok–Whitmeyer) is reported, reinforcing the second-best nature of the corrective mechanism.
\end{comment}

\chapter{Data Sources and Replication}
\label{app:data}

This appendix documents the primary data sources underlying the empirical comparisons in Chapter~\ref{ch:comparison} and the illustrative calibrations used elsewhere in the book. It also records the construction of the three comparison figures and provides guidance for readers who wish to reproduce or update the quantitative material.

\section{Primary Data Sources}

\subsection*{B.1 Taxation}

\begin{itemize}
\item Australian Taxation Office, \textit{Tax rates 2025--26} and \textit{Tax rates 2026--27}, available at \url{https://www.ato.gov.au}.  
  Personal income-tax brackets and rates for resident individuals (including the reduction of the second marginal rate to 15~per cent from 1~July 2026), company tax rates (25~per cent for base-rate entities, 30~per cent otherwise), and the Goods and Services Tax rate of 10~per cent.

\item Australian Taxation Office, \textit{Company tax rates} and associated guidance on base-rate entity eligibility (aggregated turnover threshold of \$50~million).
\end{itemize}

\subsection*{B.2 National Accounts, Prices and Labour Market}

\begin{itemize}
\item Australian Bureau of Statistics, \textit{Australian National Accounts: National Income, Expenditure and Product}, March quarter 2026 (Cat.\ No.\ 5206.0).  
  Real GDP growth (year-ended 2.5~per cent) and related aggregates.

\item Australian Bureau of Statistics, \textit{Consumer Price Index, Australia}, May 2026 (Cat.\ No.\ 6401.0).  
  Headline CPI inflation of 4.0~per cent.

\item Australian Bureau of Statistics, \textit{Labour Force, Australia}, June 2026 (Cat.\ No.\ 6202.0).  
  Unemployment rate of 4.4~per cent and associated labour-market indicators.

\item Australian Bureau of Statistics, \textit{Government Finance Statistics, Annual} (various issues, Cat.\ No.\ 5512.0).  
  Government spending as a share of GDP (approximately 28--29~per cent in the most recent observations).

\item Australian Bureau of Statistics, productivity estimates drawn from the National Accounts and the Labour Account (various releases).  
  Labour-productivity growth near zero or only marginally positive in recent years.
\end{itemize}

\subsection*{B.3 Monetary Policy and Financial Indicators}

\begin{itemize}
\item Reserve Bank of Australia, \textit{Chart Pack} (June 2026 and subsequent updates).  
  Cash-rate target, inflation measures, GDP growth, and related macroeconomic series.

\item Reserve Bank of Australia, \textit{Statement on Monetary Policy} and statistical tables (various).  
  Confirmation of the inflation target of 2--3~per cent on average over the cycle and the operational framework for the cash-rate target.
\end{itemize}

\subsection*{B.4 Public Ownership and State-Owned Enterprises}

\begin{itemize}
\item Department of Finance, ownership and governance information on Commonwealth entities and companies (including NBN Co, Australia Post, Snowy Hydro, Defence Housing Australia, Australian Rail Track Corporation and related bodies).

\item State Treasury publications on government-owned corporations (energy, water, ports and rail) in New South Wales, Queensland, Victoria and other jurisdictions.

\item OECD, \textit{Ownership and Governance of State-Owned Enterprises} (2024 and related editions), for comparative context on the size and sectoral distribution of SOEs.
\end{itemize}

\subsection*{B.5 Competition and Price Regulation}

\begin{itemize}
\item Treasury and Australian Competition and Consumer Commission materials on the Food and Grocery Code amendments and the excessive-pricing prohibition applicable from 1~July 2026 to supermarket chains with annual revenue exceeding \$30~billion.
\end{itemize}

\subsection*{B.6 Secondary and Analytical Sources}

\begin{itemize}
\item Institute of Public Affairs analysis of ABS national-accounts data on the evolving share of government spending in GDP (cited for the approximate 28--29~per cent figure).

\item Albrecht, B., Tabarrok, A.\ and Whitmeyer, M.\ (2026), ``Chaos and Misallocation under Price Controls,'' arXiv:2602.12066, for the theoretical bounds on cross-market misallocation losses used in Appendix~\ref{app:math}.
\end{itemize}

All web-based sources were current as of July 2026. Readers should consult the originating agencies for any subsequent revisions.

\section{Construction of Comparison Figures}

Three figures in Chapter~\ref{ch:comparison} were generated from the sources listed above.

\begin{description}
\item[Figure~\ref{fig:tax} (Tax Rate Comparison).]  
  Marginal and statutory rates for personal income (low, mid and high brackets), company tax (base-rate and standard) and GST under the 2026--27 Australian schedule are taken directly from ATO publications. The constrained-regime rates are the flat 10~per cent personal rate required by the binding constraints; company tax and GST are shown as zero in the figure solely to highlight that they are not prescribed by the seven constraints (the public equity claim functions as the functional counterpart to a corporate profit share).

\item[Figure~\ref{fig:ownership} (Government Ownership and Economic Role).]  
  The Australian bars reflect the limited SOE footprint (approximately 5~per cent illustrative equity share in the corporate sector), government spending of approximately 29~per cent of GDP, and a correspondingly small estimated SOE share of economic activity. The constrained-regime bars implement the mandatory 30~per cent equity claim and the associated projection of government spending and SOE-equivalent weight.

\item[Figure~\ref{fig:macro} (Macroeconomic Performance Indicators).]  
  Australian values are the observed 2026 outcomes cited in Section~B.2. Constrained-regime values are comparative-static illustrations derived from the formal model of Chapters~\ref{ch:model} and~\ref{ch:solutions}, incorporating residual efficiency losses after application of the second-best mechanism. They are not forecasts. Calibration details and sensitivity to alternative residual misallocation ratios appear in Appendix~\ref{app:math}.
\end{description}

\section{Replication Guidance}

The figures were produced with Python 3 (matplotlib, numpy) in a controlled environment. Readers wishing to reproduce or update them should:

\begin{enumerate}
\item Obtain the latest releases of the ABS and ATO series listed above.
\item Substitute the updated numerical values into the plotting scripts (available from the author on request).
\item For the constrained-regime macro column, apply the same comparative-static logic documented in Appendix~\ref{app:math}, adjusting the residual misallocation ratio $L_M/L_H$ as desired.
\end{enumerate}

No proprietary data or restricted-access microdata were used. All quantitative claims in the main text can be verified from the publicly available sources enumerated in this appendix.

\begin{comment}
NOTES FOR APPENDIX B (fully expanded into the text above):
- Full list of ABS catalogue numbers, RBA publications, ATO rate schedules, Department of Finance and state Treasury SOE information, OECD comparative data, and the Albrecht--Tabarrok--Whitmeyer reference.
- Explicit notes on how each of the three comparison figures was constructed from those sources.
- Replication guidance for readers who wish to update or re-generate the figures; no proprietary data employed.
\end{comment}

\chapter{Glossary of Institutional Terms}
\label{app:glossary}

This glossary provides short, self-contained definitions of the principal institutional and analytical terms used in the book. It is intended for readers who may not be specialists in economic theory or comparative political economy. Terms are listed alphabetically.

\begin{description}
\item[Agency cost]  
The loss in firm value that arises when the interests of managers (or other agents) diverge from the interests of residual claimants (principals). In the constrained regime the mandatory public equity claim of 30~per cent creates a specific agency cost by attenuating the manager’s residual return.

\item[Chaos Theorem (price controls)]  
The result, established by Albrecht, Tabarrok and Whitmeyer (2026), that under a binding price ceiling suppliers become indifferent across destinations, so that arbitrarily small cost differences can tip the allocation into extreme corner solutions. Some markets are fully served while others are starved, generating cross-market misallocation losses that can substantially exceed the classical Harberger triangle.

\item[Command socialism]  
An economic system characterised by near-complete state ownership of the means of production, administrative determination of prices and output targets, and comprehensive material-balance planning. Market coordination is residual or illicit.

\item[Corporatism (equity-share / state-equity)]  
An organisational principle in which the state is embedded inside firms through ownership or mandatory institutional representation rather than remaining solely an external regulator. In the present work the term refers specifically to the mandatory 30~per cent public residual claim that places the Commonwealth inside every corporate boardroom.

\item[Equity claim (residual claim)]  
The right to the net surplus of a firm after all contractual obligations (wages, interest, input costs, taxes) have been met. The residual claimant bears the residual risk and therefore has the strongest incentive to maximise the joint surplus. In the constrained regime this claim is split 30~per cent public / 70~per cent private for every corporation.

\item[Euler equation (distorted)]  
The intertemporal optimality condition that governs capital accumulation. When evaluated at frozen official prices rather than at shadow prices, it generally implies a steady-state capital stock different from (typically lower than) the efficient level.

\item[Fascist corporatism]  
The historical inter-war form of corporatism in which private nominal title to capital was preserved while firms were organised into state-supervised corporations; prices, wages and investment were heavily directed; and residual claims were subordinated to national political objectives.

\item[Flat tax]  
A personal income tax levied at a single constant rate on all taxable income. In the constrained regime the rate is fixed at 10~per cent.

\item[Harberger triangle]  
The classical deadweight-loss triangle that measures the net loss of consumer and producer surplus caused by a price distortion that reduces the quantity transacted below the efficient level.

\item[Mixed economy (social-democratic)]  
An economic system that retains predominant private ownership and market price formation while subjecting both to significant fiscal redistribution, independent monetary stabilisation, and limited sectoral regulation. Contemporary Australia is treated as a leading example.

\item[Priority rationing (productivity-priority)]  
A non-price allocation rule that ranks claimants according to a verifiable productivity score and serves higher-scoring claimants first. In the constrained regime the rule is the principal instrument for allocating goods whose official prices are frozen.

\item[Quantity-constrained equilibrium]  
An equilibrium concept applicable when official prices are fixed and therefore cannot equate notional demand and supply. Effective transactions are determined by rationing schemes; aggregate feasibility is still required.

\item[Residual claimancy]  
See Equity claim.

\item[Second-best]  
A situation in which a first-best (fully efficient) allocation is unattainable because of binding constraints. The analytical task is to minimise the residual efficiency losses while respecting those constraints. All corrective instruments developed in this book are second-best instruments.

\item[Soft budget constraint]  
The expectation by a firm that the government (or another external agent) will provide additional funds in the event of losses. The expectation weakens exit incentives and encourages excessive risk-taking or waste. The term originates with János Kornai’s analysis of socialist economies.

\item[Shadow price]  
The true marginal social value of a good or factor—the price that would prevail in an efficient allocation. When official prices are frozen, shadow prices generally diverge from official prices and are not directly observable.

\item[State-equity corporatism with frozen prices]  
The compact formal designation of the constrained regime developed in this book: a hybrid system defined by a mandatory 30~per cent public residual claim, universal freezing of official commodity prices, residual private ownership, and selective factor-price flexibility.
\end{description}

\begin{comment}
NOTES FOR APPENDIX C (fully expanded into the text above):
- Short, precise definitions of residual claimancy, quantity-constrained equilibrium, soft budget constraint, corporatism, priority rationing, Chaos Theorem, Harberger triangle, second-best, shadow price, and the other principal terms used throughout the book.
- Written expressly for the non-specialist reader (economists in other fields, political scientists, government officials).
\end{comment}

% ----------------------------------------------------------
% BIBLIOGRAPHY
% ----------------------------------------------------------
\begin{thebibliography}{99}

\bibitem{abs2026}
Australian Bureau of Statistics (various releases, 2025--2026).
\newblock \emph{Australian National Accounts: National Income, Expenditure and Product} (Cat.~No.~5206.0);
\newblock \emph{Consumer Price Index, Australia} (Cat.~No.~6401.0);
\newblock \emph{Labour Force, Australia} (Cat.~No.~6202.0);
\newblock \emph{Government Finance Statistics, Annual} (Cat.~No.~5512.0).

\bibitem{ato2026}
Australian Taxation Office (2026).
\newblock \emph{Tax Rates 2025--26 and 2026--27}; company tax rate guidance.

\bibitem{rba2026}
Reserve Bank of Australia (2026).
\newblock \emph{Chart Pack}; \emph{Statements on Monetary Policy}; statistical tables.

\bibitem{albrecht2026}
Albrecht, B., Tabarrok, A. and Whitmeyer, M. (2026).
\newblock Chaos and Misallocation under Price Controls.
\newblock \emph{arXiv:2602.12066}.

\bibitem{kornai1980}
Kornai, J. (1980).
\newblock \emph{Economics of Shortage}.
\newblock Amsterdam: North-Holland.

\bibitem{north1990}
North, D.~C. (1990).
\newblock \emph{Institutions, Institutional Change and Economic Performance}.
\newblock Cambridge: Cambridge University Press.

\bibitem{acemoglu2012}
Acemoglu, D. and Robinson, J.~A. (2012).
\newblock \emph{Why Nations Fail: The Origins of Power, Prosperity and Poverty}.
\newblock New York: Crown.

\bibitem{williamson1985}
Williamson, O.~E. (1985).
\newblock \emph{The Economic Institutions of Capitalism}.
\newblock New York: Free Press.

\bibitem{finance2026}
Department of Finance (various).
\newblock Ownership and governance information on Commonwealth entities and companies.

\bibitem{oecd2024}
OECD (2024).
\newblock \emph{Ownership and Governance of State-Owned Enterprises 2024}.
\newblock Paris: OECD Publishing.

\bibitem{treasury2026}
Australian Treasury and Australian Competition and Consumer Commission (2025--2026).
\newblock Materials on the Food and Grocery Code amendments and the excessive-pricing prohibition.

\end{thebibliography}

% ----------------------------------------------------------
% INDEX (placeholder)
% ----------------------------------------------------------
% \printindex

\end{document}
